\documentclass[12pt,a4paper]{article}

% ------------------------------------------------------------
% Pacotes
% ------------------------------------------------------------
\usepackage[T1]{fontenc}
\usepackage[utf8]{inputenc}
\usepackage[brazil]{babel}
\usepackage{lmodern}
\usepackage{microtype}
\usepackage{geometry}
\geometry{top=2.2cm,bottom=2.2cm,left=2.5cm,right=2.5cm}
\usepackage{amsmath,amssymb,mathtools,bm}
\usepackage{graphicx}
\usepackage{booktabs}
\usepackage{array}
\usepackage{siunitx}
\sisetup{output-decimal-marker={,},group-separator={.},group-minimum-digits=4,group-digits=integer}
\usepackage{xcolor}
\usepackage{float}
\usepackage{caption}
\usepackage{subcaption}
\usepackage{enumitem}
\usepackage{fancyhdr}
\usepackage{hyperref}
\usepackage{listings}
\usepackage{tikz}
\usetikzlibrary{positioning,arrows.meta}

\definecolor{azulusp}{RGB}{20,55,110}
\definecolor{cinzaclaro}{RGB}{245,247,250}
\definecolor{cinzamedio}{RGB}{90,95,105}
\definecolor{verdecodigo}{RGB}{0,115,75}

\hypersetup{
  colorlinks=true,
  linkcolor=azulusp,
  citecolor=azulusp,
  urlcolor=azulusp,
  pdftitle={Parte 1 - Identificacao de planta conhecida por LMS},
  pdfauthor={Aluno}
}

\pagestyle{fancy}
\fancyhf{}
\lhead{PSI3431 - Processamento Estatístico de Sinais}
\rhead{Parte 1 - Planta conhecida}
\cfoot{\thepage}
\setlength{\headheight}{15pt}

\newcommand{\E}{\mathbb{E}}
\newcommand{\R}{\mathbf{R}}
\newcommand{\wtrue}{\mathbf{w}^{o}}
\newcommand{\what}{\widehat{\mathbf{w}}}
\newcommand{\wt}{\widetilde{\mathbf{w}}}
\newcommand{\uvec}{\mathbf{u}}

\newenvironment{decisao}{%
  \begin{quote}\small
  \noindent\textbf{\color{azulusp}Decisão metodológica explícita.}\par\smallskip
}{%
  \end{quote}
}

\lstdefinestyle{matlabstyle}{
  language=Matlab,
  basicstyle=\ttfamily\footnotesize,
  keywordstyle=\color{azulusp}\bfseries,
  commentstyle=\color{verdecodigo},
  stringstyle=\color{purple},
  numbers=left,
  numberstyle=\tiny\color{cinzamedio},
  stepnumber=1,
  numbersep=8pt,
  frame=single,
  breaklines=true,
  showstringspaces=false,
  tabsize=4,
  backgroundcolor=\color{cinzaclaro},
  literate=
    {á}{{\'a}}1 {à}{{\`a}}1 {ã}{{\~a}}1 {â}{{\^a}}1
    {é}{{\'e}}1 {ê}{{\^e}}1 {í}{{\'i}}1
    {ó}{{\'o}}1 {ô}{{\^o}}1 {õ}{{\~o}}1
    {ú}{{\'u}}1 {ç}{{\c c}}1
    {Á}{{\'A}}1 {É}{{\'E}}1 {Í}{{\'I}}1 {Ó}{{\'O}}1 {Ú}{{\'U}}1 {Ç}{{\c C}}1
}

% ------------------------------------------------------------
% Documento
% ------------------------------------------------------------
\begin{document}

% ------------------------------------------------------------
% Capa
% ------------------------------------------------------------
\begin{titlepage}
    \centering

    {\Large Universidade de São Paulo\par}
    \vspace{0.3cm}
    {\large Escola Politécnica\par}
    \vspace{0.3cm}
    {\large PSI3431 -- Processamento Estatístico de Sinais\par}

    \vfill

    {\LARGE \textbf{Projeto Final}\par}
    \vspace{0.4cm}
    {\Large Identificação de Sistemas\par}

    \vfill

    \begin{center}
    \begin{tabular}{ll}
        \textbf{Integrante 1:} & José Victor Lyra de Castro \\
        \textbf{NUSP:} & 14747340 \\[0.4cm]

        \textbf{Integrante 2:} & Caio de Paula Colnago \\
        \textbf{NUSP:} & 12608900 \\[0.4cm]

        \textbf{Integrante 3:} & Arthur Jesus \\
        \textbf{NUSP:} & 10001000 \\
    \end{tabular}
    \end{center}

    \vfill

    {\large São Paulo\par}
    {\large Junho de 2026\par}

\end{titlepage}

% ------------------------------------------------------------
% Sumário
% ------------------------------------------------------------
\tableofcontents
\newpage

% ------------------------------------------------------------
% Introdução
% ------------------------------------------------------------
\section{Introdução}

Este relatório apresenta o desenvolvimento do projeto final da disciplina PSI3431 -- Processamento Estatístico de Sinais. O objetivo principal é realizar a identificação de sistemas por meio de algoritmos de filtragem adaptativa e estimadores baseados em erro quadrático médio.

O trabalho está dividido em duas partes. Na Parte 1, considera-se uma planta conhecida, permitindo analisar a convergência do algoritmo escolhido. Na Parte 2, a planta é desconhecida e deve ser estimada a partir de sinais de entrada e saída fornecidos.

% ------------------------------------------------------------
% Parte 1
% ------------------------------------------------------------
\section{Parte 1 -- Planta Conhecida}

\subsection{Descrição do Problema}

Nesta etapa, considera-se uma planta conhecida representada pelo vetor de coeficientes fornecido no enunciado. O objetivo é utilizar um algoritmo de filtragem adaptativa para identificar essa planta a partir de sinais de entrada brancos.

\subsection{Modelo Utilizado}

A saída da planta pode ser modelada por

\begin{equation}
    d(i) = \mathbf{u}^T(i)\mathbf{w}_o + v(i),
\end{equation}

em que $\mathbf{u}(i)$ é o vetor de entrada, $\mathbf{w}_o$ é o vetor da planta real e $v(i)$ representa o ruído.

O estimador adaptativo produz uma saída

\begin{equation}
    \hat{d}(i) = \mathbf{u}^T(i)\mathbf{w}(i),
\end{equation}

e o erro é dado por

\begin{equation}
    e(i) = d(i) - \hat{d}(i).
\end{equation}

\subsection{Algoritmo Escolhido}

O algoritmo escolhido foi o LMS, com a recursão
\begin{equation}
\boxed{
\begin{aligned}
\widehat d_m(n)&=\uvec_m^{T}(n)\mathbf{w}_m(n),\\
e_m(n)&=d_m(n)-\widehat d_m(n),\\
\mathbf{w}_m(n+1)&=\mathbf{w}_m(n)+\mu\uvec_m(n)e_m(n),
\end{aligned}}
\end{equation}
para \(m=1,\ldots,1000\), com \(\mu=0.01\), \(L=10\), \(N=1500\) e \(\mathbf{w}_m(0)=\mathbf{0}\).

A implementação em MATLAB é:
\begin{lstlisting}[style=matlabstyle,caption={Núcleo da recursão LMS em paralelo.}]
for n = 1:N
    reg(2:end, :) = reg(1:end-1, :);
    reg(1, :) = u(:, n).';

    d = w0.' * reg + v(:, n).';
    y = sum(W .* reg, 1);
    e = d - y;

    Werr = repmat(w0, 1, M) - W;
    ea = sum(Werr .* reg, 1);

    mse(n)  = mean(e.^2);
    emse(n) = mean(ea.^2);
    msd(n)  = mean(sum(Werr.^2, 1));

    W = W + mu * reg .* repmat(e, L, 1);
end
\end{lstlisting}

\subsection{Curva de MSE}

A Figura~\ref{fig:mseemse} apresenta a curva de MSE em escala logarítmica. No início, como \(\mathbf{w}(0)=\mathbf{0}\),
\begin{equation}
\mathrm{MSE}(0)\approx \sigma_u^2\|\wtrue\|^2+\sigma_v^2.
\end{equation}
Como \(\|\wtrue\|^2\approx1\), o erro parte de uma potência próxima da unidade. Após o transitório, o MSE converge para a vizinhança de \(J_{\min}=10^{-3}\), acrescida do erro em excesso causado pelo passo não nulo.

\begin{figure}[H]
\centering
\includegraphics[width=0.98\textwidth]{parte1_mse_emse.png}
\caption{Curvas médias de MSE e EMSE sobre \num{1000} realizações.}
\label{fig:mseemse}
\end{figure}

A média das 300 últimas iterações resultou em
\begin{equation}
\boxed{\mathrm{MSE}_{\infty,\,sim}\approx\num{1.0550e-3}.}
\end{equation}

\subsection{Curva de EMSE}

Para entrada branca gaussiana e LMS na convenção da Equação, a previsão de regime permanente é
\begin{equation}
\mathrm{EMSE}_{\infty}
=\frac{\mu L\sigma_u^2\sigma_v^2}
{2-\mu(L+2)\sigma_u^2}.
\label{eq:emse_theory}
\end{equation}
Substituindo os parâmetros,
\begin{equation}
\boxed{\mathrm{EMSE}_{\infty,\,teo}\approx\num{5.3238e-5}.}
\end{equation}
Consequentemente,
\begin{equation}
\boxed{\mathrm{MSE}_{\infty,\,teo}
=\sigma_v^2+\mathrm{EMSE}_{\infty,\,teo}
\approx\num{1.05324e-3}.}
\end{equation}

A simulação forneceu
\begin{equation}
\boxed{\mathrm{EMSE}_{\infty,\,sim}\approx\num{5.2808e-5},}
\end{equation}
valor muito próximo da previsão teórica. O desajuste relativo é
\begin{equation}
\mathcal{M}=\frac{\mathrm{EMSE}_{\infty}}{J_{\min}}
\approx\num{0.0532}=\num{5.32}\%.
\end{equation}

A Figura~\ref{fig:msd} apresenta o MSD. Como \(\widehat\sigma_u^2\approx1\), as equações implicam que EMSE e MSD possuem praticamente o mesmo valor numérico em regime.

\begin{figure}[H]
\centering
\includegraphics[width=0.94\textwidth]{parte1_msd.png}
\caption{Desvio quadrático médio dos coeficientes.}
\label{fig:msd}
\end{figure}

\subsection{Análise dos Resultados}

\subsubsection*{Concordância entre teoria e simulação}

\begin{table}[H]
\centering
\caption{Comparação entre regime permanente simulado e teórico.}
\begin{tabular}{lrrr}
\toprule
Grandeza & Simulação & Teoria & Erro relativo\\
\midrule
MSE & \num{1.0550e-3} & \num{1.05324e-3} & \num{0.17}\%\\
EMSE & \num{5.2808e-5} & \num{5.3238e-5} & \num{-0.81}\%\\
MSD & \num{5.3199e-5} & \num{5.3194e-5} & \num{0.01}\%\\
\bottomrule
\end{tabular}
\end{table}

A concordância é forte e sustenta as hipóteses adotadas: entrada aproximadamente branca, ruído independente e número suficiente de realizações para aproximar valores esperados.

\subsubsection*{Convergência dos coeficientes}

A Figura~\ref{fig:coeftraj} mostra a evolução da média dos dez coeficientes. Todos convergem para seus respectivos valores verdadeiros e permanecem flutuando em torno deles em regime permanente.

\begin{figure}[H]
\centering
\includegraphics[width=0.98\textwidth]{parte1_coeficientes.png}
\caption{Trajetórias médias dos coeficientes estimados; linhas tracejadas indicam os valores verdadeiros.}
\label{fig:coeftraj}
\end{figure}

A comparação final é mostrada na Figura~\ref{fig:coefcomp} e na Tabela~\ref{tab:coeficientes}.

\begin{figure}[H]
\centering
\includegraphics[width=0.92\textwidth]{comparacao_coeficientes.png}
\caption{Planta verdadeira e estimativa média ao final de \num{1500} iterações.}
\label{fig:coefcomp}
\end{figure}

\begin{table}[H]
\centering
\caption{Coeficientes verdadeiros e estimados.}
\label{tab:coeficientes}
\small
\begin{tabular}{rrrr}
\toprule
\(k\) & \(w_k^o\) & \(\overline{w}_k(N)\) & Erro médio final\\
\midrule
0 & -0.524700 & -0.524638 & \num{6.2e-5}\\
1 & -0.206000 & -0.206065 & \num{-6.5e-5}\\
2 &  0.332400 &  0.332382 & \num{-1.8e-5}\\
3 & -0.263100 & -0.262939 & \num{1.61e-4}\\
4 &  0.235800 &  0.235944 & \num{1.44e-4}\\
5 &  0.030400 &  0.030408 & \num{8.0e-6}\\
6 &  0.352500 &  0.352568 & \num{6.8e-5}\\
7 & -0.481200 & -0.481204 & \num{-4.0e-6}\\
8 & -0.048500 & -0.048565 & \num{-6.5e-5}\\
9 & -0.296400 & -0.296333 & \num{6.7e-5}\\
\bottomrule
\end{tabular}
\end{table}

A norma do erro entre o vetor médio final e a planta verdadeira é aproximadamente
\begin{equation}
\|\overline{\mathbf{w}}(N)-\wtrue\|\approx\num{2.62e-4},
\end{equation}
confirmando que o viés médio final é desprezível.

\subsubsection*{Interpretação do MSE e do EMSE}

O MSE não converge a zero porque inclui a potência do ruído irredutível. Já o EMSE mede somente a parcela causada pelo desalinhamento do filtro. A separação
\begin{equation}
\mathrm{MSE}=\sigma_v^2+\mathrm{EMSE}
\end{equation}
explica a diferença vertical entre as duas curvas após a convergência. No regime, o MSE fica próximo de \(-29.8\) dB e o EMSE próximo de \(-42.7\) dB.

\subsubsection*{Compromisso associado ao passo}

O passo \(\mu\) estabelece o compromisso entre rapidez de convergência e erro em regime. Um passo maior reduziria o tempo de convergência, porém aumentaria o EMSE e aproximaria o algoritmo do limite de instabilidade. Um passo menor reduziria o desajuste, mas exigiria mais iterações.

% ------------------------------------------------------------
% Parte 2
% ------------------------------------------------------------
\section{Parte 2 -- Planta Desconhecida}

\subsection{Descrição do Problema}

Nesta etapa, a planta é desconhecida. Foram fornecidos sinais de entrada e saída para dois cenários: entrada branca e entrada correlacionada. A partir desses dados, busca-se estimar a ordem da planta, o vetor de coeficientes e a curva de aprendizado.

\subsection{Organização dos Dados}

Os arquivos \texttt{whites.mat} e \texttt{correl.mat} contêm as variáveis \texttt{X} e \texttt{Y}, em que \texttt{X} representa os sinais de entrada e \texttt{Y} representa as respectivas saídas da planta.

\subsection{Estimativa da Ordem da Planta}

A ordem foi estimada testando diferentes comprimentos $L$ para o filtro FIR identificador. Para cada valor de $L$, calculou-se o erro quadrático médio normalizado. O valor escolhido foi aquele a partir do qual o erro deixa de apresentar redução significativa.

\begin{figure}[H]
    \centering
    \includegraphics[width=0.9\textwidth]{ordemBranco.png}
    \caption{Escolha da ordem da planta para o caso com entrada branca.}
    \label{fig:ordem-branco}
\end{figure}

\begin{figure}[H]
    \centering
    \includegraphics[width=0.9\textwidth]{ordemCorrelacionado.png}
    \caption{Escolha da ordem da planta para o caso com entrada correlacionada.}
    \label{fig:ordem-correlacionado}
\end{figure}

Pelos resultados, o número de coeficientes estimado foi $L = 31$. Portanto, como a ordem de um filtro FIR é $L-1$, a ordem estimada da planta é

\begin{equation}
    M = 30.
\end{equation}

\subsection{Curva de Aprendizado}

A curva de aprendizado foi obtida calculando o erro quadrático médio ao longo das iterações. O erro utilizado foi

\begin{equation}
    e(i) = d(i) - \hat{d}(i),
\end{equation}

e o MSE foi estimado pela média de $e^2(i)$ ao longo das realizações disponíveis.

\begin{figure}[H]
    \centering
    \includegraphics[width=0.9\textwidth]{mseAprendizadoBranco.png}
    \caption{Curva de aprendizado para o caso com entrada branca.}
    \label{fig:mse-branco}
\end{figure}

\begin{figure}[H]
    \centering
    \includegraphics[width=0.9\textwidth]{mseAprendizadoCorrelacionado.png}
    \caption{Curva de aprendizado para o caso com entrada correlacionada.}
    \label{fig:mse-correlacionado}
\end{figure}

\begin{figure}[H]
    \centering
    \includegraphics[width=0.9\textwidth]{comparacaoMSE.png}
    \caption{Comparação das curvas de aprendizado para entrada branca e correlacionada.}
    \label{fig:comparacao-mse}
\end{figure}

Observa-se que o caso com entrada branca converge mais rapidamente, enquanto o caso com entrada correlacionada apresenta convergência mais lenta. Esse comportamento é esperado, pois sinais correlacionados tornam o problema de identificação mais mal condicionado.

\subsection{Vetor Estimado da Planta}

O vetor estimado obtido para o caso de entrada branca foi

\begin{equation}
\begin{aligned}
\hat{\mathbf{w}}_{\text{branco}} = [&
0.4779,\ 0.1357,\ 0.2266,\ -0.1739,\ -0.0770,\ -0.0448,\\
&0.1804,\ -0.0458,\ 0.1153,\ -0.3372,\ -0.0581,\ -0.1354,\\
&-0.2592,\ 0.0835,\ 0.0463,\ 0.0055,\ -0.2192,\ 0.1852,\\
&0.0576,\ -0.0492,\ 0.0039,\ -0.0431,\ -0.2875,\ -0.0470,\\
&-0.1366,\ -0.1609,\ -0.1899,\ -0.0877,\ -0.3291,\\
&0.1585,\ 0.0855]^T.
\end{aligned}
\end{equation}

O vetor estimado para o caso de entrada correlacionada foi

\begin{equation}
\begin{aligned}
\hat{\mathbf{w}}_{\text{corr}} = [&
0.4779,\ 0.1356,\ 0.2266,\ -0.1740,\ -0.0769,\ -0.0448,\\
&0.1806,\ -0.0457,\ 0.1153,\ -0.3372,\ -0.0580,\ -0.1355,\\
&-0.2592,\ 0.0835,\ 0.0463,\ 0.0055,\ -0.2192,\ 0.1853,\\
&0.0575,\ -0.0491,\ 0.0038,\ -0.0430,\ -0.2876,\ -0.0470,\\
&-0.1366,\ -0.1610,\ -0.1899,\ -0.0877,\ -0.3291,\\
&0.1585,\ 0.0855]^T.
\end{aligned}
\end{equation}

\begin{figure}[H]
    \centering
    \includegraphics[width=0.9\textwidth]{vetorEstimadoBranco.png}
    \caption{Vetor estimado da planta para o caso com entrada branca.}
    \label{fig:vetor-branco}
\end{figure}

\begin{figure}[H]
    \centering
    \includegraphics[width=0.9\textwidth]{vetorEstimadoCorrelacionado.png}
    \caption{Vetor estimado da planta para o caso com entrada correlacionada.}
    \label{fig:vetor-correlacionado}
\end{figure}

Os dois vetores estimados são praticamente iguais, indicando consistência no processo de identificação da planta.

\subsection{Análise dos Resultados}

O erro final normalizado obtido foi aproximadamente $0{,}001002$ para o caso com entrada branca e $0{,}0006399$ para o caso com entrada correlacionada. Ambos os valores são baixos, indicando que o modelo estimado representa bem a planta desconhecida.

A diferença principal entre os dois cenários está na velocidade de convergência. A entrada branca permite convergência mais rápida, enquanto a entrada correlacionada exige mais iterações para atingir o mesmo patamar de erro.

% ------------------------------------------------------------
% Conclusão
% ------------------------------------------------------------
\section{Conclusão}

Neste projeto, foram aplicados conceitos de identificação de sistemas e filtragem adaptativa. Na Parte 1, a planta conhecida permitiu avaliar a convergência do algoritmo escolhido. Na Parte 2, a planta desconhecida foi estimada a partir de sinais de entrada e saída.

Os resultados da Parte 2 indicaram uma planta FIR de ordem 30, representada por 31 coeficientes. Além disso, os vetores estimados nos casos de entrada branca e correlacionada foram praticamente coincidentes, mostrando consistência na identificação.

\end{document}